\documentclass[11pt]{amsart}
\usepackage[T1]{fontenc}
\usepackage{lmodern}
\usepackage{amsmath,amssymb,mathtools}
\usepackage[margin=1in]{geometry}
\usepackage{microtype}
\usepackage{needspace}
\usepackage{hyperref}
\hypersetup{colorlinks=true,linkcolor=blue,citecolor=blue,urlcolor=blue,
  pdftitle={Asymptotic Counts of Goldbach Decompositions in the Nonnegative Polynomial Semiring},
  pdfsubject={Uniform fixed-degree counting of ordered sums of multiplicative semiring atoms}}
\urlstyle{same}

\newcommand{\N}{\mathbb N}
\newcommand{\Nzero}{\mathbb N_0}
\newcommand{\Z}{\mathbb Z}
\newcommand{\Q}{\mathbb Q}
\newcommand{\F}{\mathbb F}
\newcommand{\Prob}{\mathbb P}
\newcommand{\Ex}{\mathbb E}
\newcommand{\Atoms}{\mathcal A}
\newcommand{\Sseries}{\mathfrak S}
\newcommand{\ind}{\mathbf 1}
\DeclareMathOperator{\cont}{cont}
\DeclareMathOperator{\TV}{TV}
\DeclareMathOperator{\Law}{Law}

\newtheorem{theorem}{Theorem}[section]
\newtheorem{proposition}[theorem]{Proposition}
\newtheorem{lemma}[theorem]{Lemma}
\newtheorem{corollary}[theorem]{Corollary}
\theoremstyle{definition}
\newtheorem{example}[theorem]{Example}
\theoremstyle{remark}
\newtheorem{remark}[theorem]{Remark}

\title[Asymptotic Goldbach counts]{Asymptotic Counts of Goldbach Decompositions\\
in the Nonnegative Polynomial Semiring}
\author{Shunyao Yan}
% \address{Institution or affiliation, if applicable}
\email{shunyaoyan2009@gmail.com}
\date{Draft revised October 27, 2026}
\subjclass[2020]{Primary 11P32; Secondary 16Y60, 11N37}
\keywords{Polynomial semiring, multiplicative atom, Goldbach decomposition,
singular series, primitive polynomial, uniform asymptotic counting}

\begin{document}

\begin{abstract}
For fixed integers $d\ge1$ and $k\ge2$, we count ordered representations
of a polynomial $f=\sum_{i=0}^d c_i x^i\in\Nzero[x]$ as a sum of
exactly $k$ multiplicative atoms of $\Nzero[x]$. As the smallest
coefficient tends to infinity, we prove an asymptotic formula whose
main term is the total number of coefficientwise decompositions
multiplied by an explicit, positive Euler product. The product depends
on the target only through the prime divisors of its content. The error
is uniform in the relative sizes of the coefficients. The proof combines
residue counting in weak compositions, a large-prime estimate controlled
by the shortest coefficient row, and a finite-field comparison of
primitive and atomic summands. We also obtain a content-independent
positive lower bound for the singular series, an explicit formula for
two summands, and the average value $\zeta(d+1)^{-k}$ of the normalized
representation count over expanding coefficient boxes.
\end{abstract}

\maketitle

\section{Introduction and main result}\label{sec:introduction}

Goldbach-type questions for polynomial rings concern representations as
sums of multiplicatively irreducible polynomials; see, for example,
Pollack~\cite{Pollack}. In a semiring of nonnegative polynomials, the
absence of subtraction changes both the multiplicative structure and
the available additive representations. Gotti and Polo~\cite{GottiPolo}
studied factorization in polynomial semidomains. Liao and
Polo~\cite{LiaoPolo} proved a two-atom Goldbach theorem for nonmonomial
Laurent polynomials with coefficient sum greater than three, and Kaplan
and Polo~\cite{KaplanPolo} extended this direction to more general
Laurent polynomial and series semidomains. In an earlier unpublished work, Mittal, Pieleanu, Yan studied the asymptotic proportion of two-atom Goldbach decompositions in $\Nzero[x]$. Building on this work, we generalize the analysis to ordered decompositions into exactly $k$ multiplicative atoms for every fixed integer $k\ge2$, obtaining an explicit asymptotic formula.

Here we study an enumeration problem in the ordinary semiring
$\Nzero[x]$. We fix both the degree of the target and the number of
summands, and allow all target coefficients to grow, without requiring
comparable rates of growth. The main term records the congruence
dependence between summands whose sum is prescribed. In particular,
the individual densities of primitive coefficient vectors cannot
simply be multiplied: the content of the target remains visible in
the answer. The familiar zeta factors in primitive-vector density
\cite{Nymann} reappear after averaging over the target.

Throughout, $\N=\{1,2,\ldots\}$ and $\Nzero=\{0,1,2,\ldots\}$.
An \emph{atom} of $\Nzero[x]$ is a nonzero nonunit $h$ such that
$h=uv$, with $u,v\in\Nzero[x]$, implies that one of $u,v$ is a unit.
The only unit is $1$; write $\Atoms$ for the set of atoms.
Atomicity in this paper is always multiplicative atomicity in
$\Nzero[x]$, not irreducibility over $\Q$. For example, $1+x^3$ is
an atom of $\Nzero[x]$, although it factors in $\Q[x]$.

For $k\ge2$, define
\[
G_k(f)=\#\bigl\{(g_1,\ldots,g_k)\in\Atoms^k:
                       g_1+\cdots+g_k=f\bigr\}.
\]
The summands are ordered, repetitions are allowed, and their degrees
need not equal the degree of $f$. Prime constants and the atom $x$ are
included. For
\begin{equation}\label{eq:target}
f(x)=\sum_{i=0}^d c_i x^i,\qquad c_i\in\N,
\end{equation}
put
\begin{equation}\label{eq:notation}
n=d+1,\qquad C=\gcd(c_0,\ldots,c_d),\qquad
M=\min_{0\le i\le d}(c_i+1).
\end{equation}
The number of all ordered coefficientwise decompositions, without any
atomicity condition, is
\begin{equation}\label{eq:volume}
V_k(f)=\prod_{i=0}^d\binom{c_i+k-1}{k-1}.
\end{equation}
Here the notation $V_k$ denotes a count, not a continuous volume.

For a prime $p$, set
\begin{equation}\label{eq:rho}
\rho_{p,n,k}(C)
=(1-p^{-n})^k
 +(-1)^k\bigl(p^n\ind_{p\mid C}-1\bigr)p^{-nk},
\end{equation}
and define the singular series
\begin{equation}\label{eq:series}
\Sseries_{n,k}(C)=\prod_p\rho_{p,n,k}(C).
\end{equation}

\begin{theorem}[Uniform fixed-degree, fixed-length asymptotic]
\label{thm:main}
Fix $d\ge1$ and $k\ge2$, and let $f$ satisfy \eqref{eq:target}.
The product \eqref{eq:series} converges absolutely and satisfies
\begin{equation}\label{eq:positive}
\frac{3}{4\zeta(d+1)^k}\le\Sseries_{d+1,k}(C)\le1.
\end{equation}
As $M\to\infty$,
\begin{equation}\label{eq:main}
G_k(f)=
\bigl(\Sseries_{d+1,k}(C)+o_{d,k}(1)\bigr)
\prod_{i=0}^d\binom{c_i+k-1}{k-1}.
\end{equation}
The error is uniform over all positive coefficient vectors: for every
$\varepsilon>0$ there exists $M_0(d,k,\varepsilon)$ such that
\[
\left|\frac{G_k(f)}{V_k(f)}-\Sseries_{d+1,k}(C)\right|
\le\varepsilon
\quad\text{whenever }M\ge M_0(d,k,\varepsilon).
\]
Consequently, \eqref{eq:main} is also a uniform relative asymptotic,
and every target with sufficiently large $M$ is a sum of exactly $k$
atoms.
\end{theorem}

The parameters $d$ and $k$ are fixed in this theorem. No bound on
$\max_i c_i/\min_i c_i$ is imposed. In particular, the content $C$
may vary with the target and need not tend to a limit. We make no
claim here for growing $d$ or $k$, or for targets with fixed zero
coefficients.

The proof has two parts. First we count decompositions with primitive
coefficient columns and derive their local factors by
inclusion--exclusion. The large-prime tail is controlled using the
shortest row of coefficients. Second, we show that primitive and atomic
columns disagree with probability tending uniformly to zero. For that
step, we prove the needed irreducibility-density estimate directly for
the nonuniform column distributions arising from decompositions.

\section{Coefficient distributions and residues}\label{sec:residues}

For a nonnegative integer $c$, let
\[
\mathcal C_k(c)=\{(a_1,\ldots,a_k)\in\Nzero^k:a_1+\cdots+a_k=c\}.
\]
Choose a vector uniformly from $\mathcal C_k(c_i)$, independently for
each $i=0,\ldots,d$. Writing its entries as $A_{i,1},\ldots,A_{i,k}$
and setting
\[
g_j(x)=\sum_{i=0}^d A_{i,j}x^i
\]
gives the uniform distribution on the $V_k(f)$ decompositions in
\eqref{eq:volume}. Rows are independent; columns generally are not.

The content of a nonzero integer polynomial is the nonnegative gcd of
its coefficients, and we set $\cont(0)=0$. A polynomial or coefficient
column is \emph{primitive} when its content is one. Let $P_k(f)$ count
the decompositions whose $k$ columns are all primitive. Thus
$P_k(f)/V_k(f)$ and $G_k(f)/V_k(f)$ are probabilities in the preceding
model.

For probability measures on a finite set, write
\[
\TV(\mu,\nu)=\frac12\sum_a|\mu(a)-\nu(a)|.
\]
The discrepancy on any event is at most this quantity. Marginalization
cannot increase it, and for product measures it is at most the sum of
the discrepancies of the individual factors. The latter assertion
follows by telescoping the difference of the two products.

\begin{lemma}[Residues of a weak composition]\label{lem:row}
Let $q\ge2$, and choose a uniform vector from $\mathcal C_k(c)$.
The compatible residue vectors modulo $q$ are those whose coordinate
sum is $c$ modulo $q$, and there are $q^{k-1}$ of them. If $c\ge2kq$,
the residue law has total-variation distance at most
\begin{equation}\label{eq:rowTV}
L_k\frac qc,\qquad
L_k=\frac{k(k-1)}2\left(\frac32\right)^{k-2},
\end{equation}
from the uniform law on these compatible vectors.
\end{lemma}

\begin{proof}
Fix representatives $r_1,\ldots,r_k\in\{0,\ldots,q-1\}$ with compatible
sum, and write $t=\sum_jr_j$. The entries of a lift are
$r_j+qb_j$, where
\[
b_1+\cdots+b_k=(c-t)/q.
\]
The exact number of lifts is
\begin{equation}\label{eq:lifts}
\binom{(c-t)/q+k-1}{k-1},
\end{equation}
interpreted as zero when $t>c$. Under $c\ge2kq$ we have $t<c$.
The ratio of the probability of this residue vector to $q^{1-k}$ is
\[
\prod_{h=1}^{k-1}\frac{c-t+qh}{c+h}.
\]
Each factor differs from one in absolute value by at most $kq/c\le1/2$.
Telescoping the product gives a difference from one of at most
\[
(k-1)\left(\frac32\right)^{k-2}\frac{kq}{c}.
\]
Averaging this bound over the $q^{k-1}$ residues and dividing by two
proves \eqref{eq:rowTV}.
\end{proof}

Set
\begin{equation}\label{eq:Ank}
A_{n,k}=nL_k.
\end{equation}

\begin{corollary}[Matrix and column residues]\label{cor:matrix}
If $M-1\ge2kq$, the full coefficient matrix modulo $q$ differs in total
variation by at most
\begin{equation}\label{eq:matrixTV}
A_{n,k}\frac{q}{M-1}
\end{equation}
from the uniform law on matrices with the prescribed row sums modulo
$q$. Under this comparison law, each designated column is uniform on
$(\Z/q\Z)^n$. The same bound therefore holds for its marginal law.
\end{corollary}

\begin{proof}
Apply Lemma~\ref{lem:row} to each row and add the bounds. In a
compatible row, fixing one entry leaves $q^{k-2}$ choices for the
remaining entries, including the case $k=2$. Thus that entry is
uniform modulo $q$. Independence of the rows proves the column
assertion.
\end{proof}

\begin{lemma}[A decreasing entry distribution]\label{lem:decreasing}
A designated entry in a uniform vector from $\mathcal C_k(c)$ has mass
function
\begin{equation}\label{eq:mass}
w_c(a)=\frac{\binom{c-a+k-2}{k-2}}{\binom{c+k-1}{k-1}},
\qquad0\le a\le c.
\end{equation}
It is nonincreasing, and
\[
w_c(0)=\frac{k-1}{c+k-1}.
\]
For every positive integer $p$,
\begin{equation}\label{eq:divisibility}
\Prob(p\mid A_{i,j})\le w_{c_i}(0)+\frac1p.
\end{equation}
In particular, if
\begin{equation}\label{eq:beta}
\beta_M=\frac{k-1}{M+k-2},
\end{equation}
then $w_{c_i}(0)\le\beta_M$ for every row, and
\begin{equation}\label{eq:kp}
\Prob(p\mid A_{i,j})\le\frac{k}{p}
\quad\text{for }p\le M-1.
\end{equation}
\end{lemma}

\begin{proof}
Fixing the designated entry leaves a weak composition of $c-a$ into
$k-1$ parts, which gives \eqref{eq:mass}. The binomial coefficient in
the numerator is nonincreasing in $a$.

For any nonincreasing probability mass $w$ on $\Nzero$, extended by zero
outside its support, the disjoint blocks
$\{(m-1)p+1,\ldots,mp\}$ give
\[
p\sum_{m\ge1}w(mp)\le\sum_{a\ge1}w(a)\le1.
\]
Separate the zero entry to obtain \eqref{eq:divisibility}. Finally,
$\beta_M\le(k-1)/(M-1)\le(k-1)/p$ when $p\le M-1$.
\end{proof}

\section{The primitive count and its singular series}\label{sec:primitive}

\begin{lemma}[The local probability]\label{lem:local}
Under uniform compatible residues modulo a prime $p$, the probability
that none of the $k$ columns is zero is $\rho_{p,n,k}(C)$ from
\eqref{eq:rho}.
\end{lemma}

\begin{proof}
Put $Q=p^n$ and $\mathbf c=(c_0,\ldots,c_d)\bmod p$. There are
$Q^{k-1}$ tuples of vectors in $\F_p^n$ with sum $\mathbf c$.
Requiring any specified $r<k$ columns to vanish has probability
$Q^{-r}$. Requiring all $k$ to vanish has probability
$\ind_{\mathbf c=0}Q^{1-k}$. Inclusion--exclusion therefore gives
\[
\sum_{r=0}^{k-1}(-1)^r\binom kr Q^{-r}
+(-1)^k\ind_{\mathbf c=0}Q^{1-k}
=(1-Q^{-1})^k+(-1)^k(Q\ind_{\mathbf c=0}-1)Q^{-k}.
\]
The condition $\mathbf c=0$ is exactly $p\mid C$.
\end{proof}

For an integer $y\ge2$, write
\begin{equation}\label{eq:cutoffs}
q_y=\prod_{p\le y}p,
\qquad
T_n(y)=\sum_{m>y}m^{-n}
\le\frac{y^{1-n}}{n-1}.
\end{equation}
The final inequality uses $n\ge2$.

\begin{proposition}[Convergence and uniform positivity]\label{prop:product}
For all $n,k\ge2$ and $C\ge1$, the series \eqref{eq:series} converges
absolutely, satisfies \eqref{eq:positive}, and obeys
\begin{equation}\label{eq:producttail}
0\le\prod_{p\le y}\rho_{p,n,k}(C)-\Sseries_{n,k}(C)
\le kT_n(y).
\end{equation}
\end{proposition}

\begin{proof}
Each local factor is a probability. A union bound over the $k$ zero-column
events gives
\[
0\le1-\rho_{p,n,k}(C)\le kp^{-n}.
\]
The sum of the right side converges. Telescoping finite products of
numbers in $[0,1]$, followed by passage to the limit, proves
\eqref{eq:producttail}.

For the positive lower bound, write
\[
\rho_{p,n,k}(C)=(1-p^{-n})^k(1+e_p),
\qquad
e_p=\frac{(-1)^k(p^n\ind_{p\mid C}-1)}{(p^n-1)^k}.
\]
If $k$ is even, the only negative case is $p\nmid C$, with
$|e_p|=(p^n-1)^{-k}$. If $k$ is odd, then $k\ge3$ and the only
negative case is $p\mid C$, with $|e_p|=(p^n-1)^{1-k}$.
In either case, the negative part $u_p=\max(0,-e_p)$ satisfies
\[
u_p\le(p^n-1)^{-2}.
\]
Consequently,
\begin{equation}\label{eq:quarter}
\sum_pu_p
\le\sum_p\frac1{(p^2-1)^2}
\le\frac13\sum_{m=2}^{\infty}\frac1{m^2-1}
=\frac14.
\end{equation}
For the last equality, use
$1/(m^2-1)=\tfrac12((m-1)^{-1}-(m+1)^{-1})$.

For every finite prime set,
$\prod_p(1+e_p)\ge\prod_p(1-u_p)\ge1-\sum_pu_p\ge3/4$.
The Euler product for $\zeta(n)$ now gives
$\Sseries_{n,k}(C)\ge3/(4\zeta(n)^k)$.
This also proves that no local factor or limiting product vanishes.
\end{proof}

\begin{lemma}[Uniform removal of large common primes]\label{lem:tail}
The probability that some column has a common prime divisor greater
than $y$ is at most
\begin{equation}\label{eq:largetail}
k\beta_M+k^{n+1}T_n(y).
\end{equation}
Here a zero column is regarded as divisible by every prime.
\end{lemma}

\begin{proof}
Choose a row of minimum total $M-1$. For one designated column, its entry
in that row is zero with probability at most $\beta_M$. Otherwise any
prime dividing the entire column is at most $M-1$.

For a prime $p$ in that range, the entries of the designated column
are independent across rows. By Lemma~\ref{lem:decreasing}, the
probability that they are all divisible by $p$ is at most $(k/p)^n$.
Summing over $y<p\le M-1$ and then over the $k$ columns gives
\eqref{eq:largetail}. The bound contains no largest-coefficient term.
\end{proof}

\begin{proposition}[Primitive decompositions]\label{prop:primitive}
If $M-1\ge2kq_y$, then
\begin{equation}\label{eq:primitiveerror}
\left|\frac{P_k(f)}{V_k(f)}-\Sseries_{n,k}(C)\right|
\le
A_{n,k}\frac{q_y}{M-1}
+k\beta_M+(k^{n+1}+k)T_n(y).
\end{equation}
\end{proposition}

\begin{proof}
Let $\mathcal E_y$ be the event that no column is zero modulo any
prime at most $y$. The Chinese remainder theorem identifies the
uniform compatible law modulo $q_y$ with the product of its
compatible laws at those primes. Lemma~\ref{lem:local} and
Corollary~\ref{cor:matrix} give
\[
\left|\Prob(\mathcal E_y)-\prod_{p\le y}\rho_{p,n,k}(C)\right|
\le A_{n,k}\frac{q_y}{M-1}.
\]
Primitivity of every column implies $\mathcal E_y$. Conversely, a
column that passes this finite sieve but is not primitive has a common
prime greater than $y$. Zero columns are also covered by
Lemma~\ref{lem:tail}. Apply that lemma and then
\eqref{eq:producttail}.
\end{proof}

\section{From primitive columns to atomic summands}\label{sec:atoms}

\begin{lemma}[Atomicity and content]\label{lem:atomfacts}
The constant atoms of $\Nzero[x]$ are the positive integer primes, and
the only atom with zero constant coefficient is $x$. Every
nonconstant atom is primitive. A primitive nonconstant polynomial
that is irreducible over $\Q$ is an atom of $\Nzero[x]$.
\end{lemma}

\begin{proof}
Degree comparison shows that an invertible polynomial must be constant,
and the only unit in $\Nzero$ is $1$. A factorization of a positive
constant likewise has constant factors, proving the prime assertion.
A nonzero zero-constant polynomial is $xh$; this is atomic only when
$h=1$, and $x$ is indeed an atom.

If a nonconstant polynomial has content $c>1$, factoring out $c$
is a product of two nonunits. Thus a nonconstant atom is primitive.
For a primitive nonconstant polynomial, a constant factor in
$\Nzero[x]$ must be $1$. Any nontrivial factorization would therefore
have two positive-degree factors, contradicting rational
irreducibility.
\end{proof}

For instance, $2+2x=2(1+x)$ is irreducible as a degree-one polynomial
over $\Q$, but is not a semiring atom. Conversely, $1+x^3$ is an
atom: in a nonnegative factorization its positive factor exponents
would all have to occur in $\{0,3\}$, which precludes two positive-degree
factors. These examples explain why primitivity is needed in the
comparison.

The following elementary count provides the only finite-field
irreducibility input. We include its proof; the standard finite-field
construction and subfield facts are described in Conrad~\cite{Conrad}.

\begin{lemma}[A positive density of modular certificates]\label{lem:finitefield}
Fix $d\ge2$. If $p\ge4d^2$ is prime, the proportion of vectors in
$\F_p^{d+1}$ representing an irreducible polynomial of degree exactly
$d$ is at least
\begin{equation}\label{eq:alphad}
\alpha_d=\frac1{4d}.
\end{equation}
\end{lemma}

\begin{proof}
In a splitting field over $\F_p$, the roots of $X^{p^d}-X$ form a field:
Frobenius gives closure under addition and multiplication, and a
nonzero root has its inverse in the same set. The derivative is $-1$,
so this field has exactly $p^d$ elements.

An element's degree $r$ over $\F_p$ divides $d$. If $r<d$, it is a
root of $X^{p^r}-X$, and $r\le d/2$. Thus at most
\[
\sum_{\substack{r\mid d\\r<d}}p^r\le d p^{d/2}\le\tfrac12p^d
\]
elements have degree smaller than $d$, where the final inequality
holds for the stated range of $p$.

Every monic irreducible of degree $d$ accounts for at most $d$ elements.
There are therefore at least $p^d/(2d)$ such polynomials. Their $p-1$
nonzero scalar multiples are all distinct, so their proportion among
$p^{d+1}$ coefficient vectors is at least
\[
\frac{(p-1)p^d}{2d\,p^{d+1}}
=\frac{1-1/p}{2d}\ge\frac1{4d}.
\]
\end{proof}

\begin{lemma}[Rational reducibility under the column law]
\label{lem:Qreducible}
Fix $d\ge2$, and choose distinct primes
$p_1,\ldots,p_\ell\ge4d^2$. Put
\[
R_\ell=p_1\cdots p_\ell.
\]
If $M-1\ge2kR_\ell$, then a designated summand column satisfies
\begin{equation}\label{eq:Qerror}
\Prob\bigl(\deg g_j=d,\ g_j\text{ reducible over }\Q\bigr)
\le(1-\alpha_d)^\ell
+A_{n,k}\frac{R_\ell}{M-1}.
\end{equation}
\end{lemma}

\begin{proof}
A rationally reducible integer polynomial of degree $d$ has, by Gauss's
lemma, a factorization into two positive-degree integer polynomials.
If its reduction has degree $d$, the leading coefficients of both
factors remain nonzero, so both factor degrees are preserved. Its
reduction cannot then be an irreducible of degree $d$.

Under uniform column residues modulo $R_\ell$, CRT makes the reductions
at the $\ell$ primes independent. By Lemma~\ref{lem:finitefield},
the probability that all $\ell$ degree-preserving irreducibility
certificates fail is at most $(1-\alpha_d)^\ell$. Transfer this event
bound to the actual column law using Corollary~\ref{cor:matrix}.
\end{proof}

\begin{remark}[Degree preservation is essential]\label{rem:degree}
The polynomial
\[
(1+2x)(1+x)=1+3x+2x^2
\]
is rationally reducible, but reduces modulo two to the irreducible
linear polynomial $1+x$. Lemma~\ref{lem:Qreducible} uses irreducible
reductions of degree \emph{exactly} $d$ and therefore does not make
this erroneous inference.
\end{remark}

\begin{proposition}[Uniform primitive-to-atom comparison]\label{prop:comparison}
For $d\ge2$ and the choices in Lemma~\ref{lem:Qreducible},
\begin{equation}\label{eq:atomerror}
\left|\frac{G_k(f)-P_k(f)}{V_k(f)}\right|
\le2k\beta_M+k(1-\alpha_d)^\ell
+kA_{n,k}\frac{R_\ell}{M-1}.
\end{equation}
For $d=1$, the left side is at most $k\beta_M$.
Consequently, for every fixed $d\ge1$ and $k\ge2$,
\begin{equation}\label{eq:atomprimitive}
\frac{G_k(f)-P_k(f)}{V_k(f)}=o_{d,k}(1)
\end{equation}
uniformly as $M\to\infty$.
\end{proposition}

\begin{proof}
For one column, discard the events that its constant or degree-$d$
coefficient is zero. Their combined probability is at most
$2\beta_M$ by Lemma~\ref{lem:decreasing}. Outside this boundary, the
column has positive constant coefficient and degree $d$.

By Lemma~\ref{lem:atomfacts}, every atomic such column is primitive;
if it is primitive but not atomic, it is rationally reducible.
Hence the probability that its primitive and atomic indicators differ
is bounded by $2\beta_M$ plus \eqref{eq:Qerror}. A union bound over
columns bounds the symmetric difference between the events
``all primitive'' and ``all atomic,'' proving \eqref{eq:atomerror}.
No independence between columns is used.

When $d=1$, every nonconstant column is atomic exactly when primitive:
a linear factorization can have only one positive-degree factor, so
the other factor divides the content. This includes $x$. Only constant
columns can cause a discrepancy, and each has probability at most
$\beta_M$ because its linear coefficient is zero.

For $d\ge2$, first fix $\ell$ large enough to make the exponential term
small, and then let $M$ grow. The primes, and hence $R_\ell$, stay fixed
in this limit. This proves the asserted uniformity.
\end{proof}

\begin{remark}[The two counting events need not contain one another]
Prime constants are atoms but not primitive, whereas the constant
polynomial $1$ is primitive but not an atom. Accordingly, one cannot
replace the symmetric-difference argument above by an unconditional
inclusion between the two sets of decompositions. For example, direct
linear counting gives
\[
\begin{array}{c|rrr}
f&G_2(f)&P_2(f)&V_2(f)\\\hline
2+30x&31&33&93\\
30+2x&49&33&93
\end{array}
\]
These finite boundary effects are included in the error term.
\end{remark}

\Needspace{12\baselineskip}
\section{Completion of the asymptotic formula}\label{sec:proof}

The preceding estimates make the order of limits explicit. With the
notation already defined, if $d\ge2$ and
\[
M-1\ge2k\max(q_y,R_\ell),
\]
Propositions~\ref{prop:primitive} and \ref{prop:comparison} give
\begin{equation}\label{eq:completeerror}
\begin{aligned}
\left|\frac{G_k(f)}{V_k(f)}-\Sseries_{n,k}(C)\right|
\le{}&3k\beta_M+(k^{n+1}+k)T_n(y)\\
&+k\left(1-\frac1{4d}\right)^\ell
+A_{n,k}\frac{q_y+kR_\ell}{M-1}.
\end{aligned}
\end{equation}
For $d=1$, the sharper bound is
\begin{equation}\label{eq:linearerror}
\left|\frac{G_k(f)}{V_k(f)}-\Sseries_{2,k}(C)\right|
\le2k\beta_M+(k^3+k)T_2(y)
+A_{2,k}\frac{q_y}{M-1},
\end{equation}
provided $M-1\ge2kq_y$.

\begin{proof}[Proof of Theorem~\ref{thm:main}]
Convergence and \eqref{eq:positive} were proved in
Proposition~\ref{prop:product}.
Fix $d,k$ and $\varepsilon>0$. First choose the prime cutoff $y$ so
that the $T_n(y)$ term in \eqref{eq:completeerror} is small. For
$d\ge2$, next choose $\ell$ and then a fixed set of $\ell$ sufficiently
large distinct primes so that the exponential term is small.
Finally choose $M$ large enough for the displayed size hypotheses and
for all remaining terms to be small. None of these choices depends
on the largest coefficient, on the aspect ratios, or on the target
content. For $d=1$, use \eqref{eq:linearerror} without the second
prime set.

This proves the uniform absolute asymptotic \eqref{eq:main}. Since
$\Sseries_{n,k}(C)\ge3/(4\zeta(n)^k)$ for every content, division by
the main-term factor gives the uniform relative asymptotic. Taking
the absolute error below half of this lower bound gives
\[
\frac{G_k(f)}{V_k(f)}\ge\frac{3}{8\zeta(n)^k}>0
\]
for all sufficiently large $M$, proving existence of $k$-atom
representations.
\end{proof}

The use of fixed prime cutoffs is deliberate. Finite-modulus
equidistribution alone does not justify allowing the modulus to grow
arbitrarily with the coefficient lengths. Here the large-prime error
is removed separately, through the shortest-row estimate, before any
claim of uniform asymptotics is made.

\section{Content dependence and averaging}\label{sec:consequences}

\begin{corollary}[Two atomic summands]\label{cor:two}
For fixed $d\ge1$ and targets with all coefficients tending to infinity,
\begin{equation}\label{eq:two}
G_2(f)=\left(
\prod_{p\mid C}(1-p^{-(d+1)})
\prod_{p\nmid C}(1-2p^{-(d+1)})+o_d(1)
\right)\prod_{i=0}^d(c_i+1),
\end{equation}
uniformly in the relative coefficient sizes.
\end{corollary}

\begin{proof}
For $k=2$, formula \eqref{eq:rho} equals $1-p^{-n}$ when $p\mid C$
and $1-2p^{-n}$ otherwise. Apply Theorem~\ref{thm:main}.
\end{proof}

The singular series depends only on the prime divisors of $C$, not on
their exponents. Its dependence on a newly introduced content prime
also changes with the parity of $k$:
\begin{equation}\label{eq:parityeffect}
\rho_{p,n,k}(p\mid C)-\rho_{p,n,k}(p\nmid C)
=(-1)^kp^{-n(k-1)}.
\end{equation}
For instance, with $u=p^{-n}$, the three-summand local factors are
$1-3u+2u^2$ when $p\mid C$ and $1-3u+3u^2$ otherwise.

\begin{example}[Two different target limits]\label{ex:content}
Fix $d\ge1$ and put $n=d+1$. For
\[
f_H(x)=\sum_{i=0}^d(H+i)x^i,
\]
the content is one, because the first two coefficients are consecutive.
Hence
\[
\frac{G_2(f_H)}{V_2(f_H)}\longrightarrow
\prod_p(1-2p^{-n}).
\]
For the targets
\[
h_H(x)=H!\,(1+x+\cdots+x^d),
\]
every prime at most $H$ divides the content. Formula \eqref{eq:two}
and the uniform product-tail estimate give
\[
\frac{G_2(h_H)}{V_2(h_H)}\longrightarrow
\prod_p(1-p^{-n})=\zeta(n)^{-1}.
\]
These are distinct limits, since each factor $1-2p^{-n}$ is strictly
smaller than $1-p^{-n}$. Thus the normalized count has no single
content-independent limit for arbitrary target sequences.
\end{example}

Averaging the target restores a simpler value. The next result averages
the \emph{normalized count for each target}; it does not replace this
average by a ratio of two separately summed counts.

\begin{theorem}[Average over target boxes]\label{thm:average}
Fix $d\ge1$, $k\ge2$, and let $c_i$ be independent and uniform on
$\{1,\ldots,H_i\}$. Put $f=\sum_i c_ix^i$ and $C=\gcd_i c_i$.
As $\min_iH_i\to\infty$,
\begin{equation}\label{eq:averages}
\Ex\,\Sseries_{d+1,k}(C)\longrightarrow\zeta(d+1)^{-k},
\qquad
\Ex\,\frac{G_k(f)}{V_k(f)}\longrightarrow\zeta(d+1)^{-k}.
\end{equation}
No restriction on the ratios of the $H_i$ is required.
\end{theorem}

\begin{proof}
Fix $y\ge2$. As every $H_i$ grows, the target coefficient vector
becomes uniform modulo $q_y$. CRT makes its reductions at the primes
$p\le y$ independent in the limiting law. At a single prime,
\[
\Prob(p\mid C)\longrightarrow p^{-n}.
\]
Taking the expectation of \eqref{eq:rho}, the correction term vanishes:
\[
\Ex\,\rho_{p,n,k}(C)\longrightarrow
(1-p^{-n})^k,
\]
because $p^n\cdot p^{-n}-1=0$. Independence in the limiting finite
residue law therefore gives
\[
\Ex\prod_{p\le y}\rho_{p,n,k}(C)
\longrightarrow\prod_{p\le y}(1-p^{-n})^k.
\]
The uniform bound \eqref{eq:producttail} permits $y\to\infty$,
proving the first limit in \eqref{eq:averages}.

For the second, fix $\varepsilon>0$ and take the threshold $M_0$ in
Theorem~\ref{thm:main}. The probability that
$\min_i(c_i+1)<M_0$ is at most
\[
\sum_{i=0}^d\frac{M_0-1}{H_i},
\]
which tends to zero. Outside this event the difference between the
normalized count and its singular series is at most $\varepsilon$;
on the event it is at most one, since both quantities lie in $[0,1]$.
Their expectations therefore differ by a quantity tending to zero.
\end{proof}

The main term thus separates two effects: the product of binomial
coefficients counts nonnegative coefficient allocations, while the
Euler product records simultaneous primitivity under the prescribed
sum. Rational reducibility and exceptional constant or degree-drop
summands contribute only to the vanishing error in the fixed-$d$,
fixed-$k$ limit. This gives both a uniform count for each target and
an arithmetic explanation of the simpler average density.

\section*{Acknowledgment}
AI tools like GPT-5.5, GPT-5.6, GPT-6 assisted mathematical exploration, proof
checking, and paper preparation.

\begin{thebibliography}{9}

\bibitem{Conrad}
K. Conrad,
\emph{Finite fields}, expository notes.
\url{https://kconrad.math.uconn.edu/blurbs/galoistheory/finitefields.pdf}.

\bibitem{GottiPolo}
F. Gotti and H. Polo,
\emph{On the arithmetic of polynomial semidomains},
Forum Math. \textbf{35} (2023), no.~5, 1179--1197.
\href{https://doi.org/10.1515/forum-2022-0091}{doi:10.1515/forum-2022-0091}.

\bibitem{KaplanPolo}
N. Kaplan and H. Polo,
\emph{A Goldbach theorem for Laurent series semidomains},
New York J. Math. \textbf{31} (2025), 1409--1426.
\url{https://nyjm.albany.edu/j/2025/31-53v.pdf}.

\bibitem{LiaoPolo}
S. Liao and H. Polo,
\emph{A Goldbach theorem for Laurent polynomials with positive integer coefficients},
Amer. Math. Monthly \textbf{131} (2024), no.~8, 704--711.
\href{https://doi.org/10.1080/00029890.2024.2363725}{doi:10.1080/00029890.2024.2363725}.

\bibitem{Nymann}
J. E. Nymann,
\emph{On the probability that $k$ positive integers are relatively prime},
J. Number Theory \textbf{4} (1972), no.~5, 469--473.
\href{https://doi.org/10.1016/0022-314X(72)90038-8}{doi:10.1016/0022-314X(72)90038-8}.

\bibitem{Pollack}
P. Pollack,
\emph{On polynomial rings with a Goldbach property},
Amer. Math. Monthly \textbf{118} (2011), no.~1, 71--77.
\href{https://doi.org/10.4169/amer.math.monthly.118.01.071}{doi:10.4169/amer.math.monthly.118.01.071}.

\end{thebibliography}
\end{document}
